\documentclass[a4paper,11pt]{article}
\usepackage{amsmath,amssymb,amsthm, tikz,titlesec,hyperref,esint,braket, graphicx, mathrsfs, enumitem}
\usepackage[a4paper,margin=2cm]{geometry}
\linespread{1.3}
\newtheorem{claim}{Claim}[section]

\newcommand\name{Park Chanwoo}   % Name of the student
\newcommand\university{KAIST} % Name of the university
\newcommand\department{Physics} % Name of the department
\newcommand\studentid{20230297} % Student ID
\newcommand\s{\,\;}


\newenvironment{solution}[1]
  {\renewcommand\qedsymbol{$\square$}\begin{proof}[\textbf{Solution#1}]}
  {\end{proof}}
\newenvironment{note}
  {\renewcommand\qedsymbol{$\blacksquare$}\begin{proof}[\textnormal{\textbf{note}}]}
  {\end{proof}}

\title{KAIST\\2026 Special Topics in Physics Quantum Electronic Devices I: Interferometers and Fractional Particles\\
Homework 3\bigskip}
\author{\textbf{\Large \name} \\
% University: \university\\
Department: \department\\
Student ID: \studentid}
\date{\today}

\begin{document}
\thispagestyle{empty}
\maketitle
\tableofcontents
\titleformat{\section}[frame]{\pagebreak}{\filright
\footnotesize  \enspace \textsf{KAIST --- Interferometers and Fractional Particles 2026 Spring}\enspace}{6pt}{\Large\bfseries\filcenter}
\setcounter{secnumdepth}{1}
\setcounter{tocdepth}{2}

\newcommand{\tr}{\operatorname{tr}}
\newcommand{\sgn}{\operatorname{sgn}}
\newcommand{\supp}{\operatorname{supp}}
\renewcommand{\Re}{\operatorname{Re}}
\renewcommand{\Im}{\operatorname{Im}}
\newcommand{\boltz}{k_{\mathrm{B}}}

\section{\#1}
\subsection{(1a)}

The boundary conditions for $V_1(x) = \alpha\delta(x)$ are
\begin{equation}
    \psi(0^+) = \psi(0^-), \qquad \psi'(0^+) - \psi'(0^-) = \frac{2m\alpha}{\hbar^2}\psi(0).
\end{equation}
With $\psi = e^{ikx} + re^{-ikx}$ for $x < 0$ and $\psi = te^{ikx}$ for $x > 0$, these become
\begin{equation}
    1 + r = t, \qquad ik(t - 1 + r) = \frac{2m\alpha}{\hbar^2}\,t.
\end{equation}
Defining $\gamma \equiv m\alpha/(\hbar^2 k)$ and substituting $r = t - 1$:
\begin{equation}
    i(2t - 2) = 2\gamma\, t \implies t(i - \gamma) = i \implies t = \frac{1}{1 + i\gamma}, \quad r = \frac{-i\gamma}{1 + i\gamma}.
\end{equation}
Therefore,
\begin{equation}
    T_1 = \frac{1}{1 + \gamma^2}, \qquad R_1 = \frac{\gamma^2}{1 + \gamma^2}.
\end{equation}


\subsection{(1b)}

For $V_2(x) = \alpha[\delta(x) + \delta(x-d)]$, write
\begin{equation}
    \psi = \begin{cases}
        e^{ikx} + re^{-ikx}, & x < 0, \\
        Ae^{ikx} + Be^{-ikx}, & 0 < x < d, \\
        te^{ikx}, & x > d.
    \end{cases}
\end{equation}

\noindent\textbf{Boundary conditions at $x = 0$:}
\begin{align}
    1 + r &= A + B, \label{eq:c1}\\
    i(A - B - 1 + r) &= 2\gamma(A + B). \label{eq:d1}
\end{align}
Substituting $r = A + B - 1$ from \eqref{eq:c1} into \eqref{eq:d1}:
\begin{equation}
    i(A - 1) = \gamma(A + B). \label{eq:AB1}
\end{equation}

\noindent\textbf{Boundary conditions at $x = d$:}
\begin{align}
    Ae^{ikd} + Be^{-ikd} &= te^{ikd}, \\
    ik(te^{ikd} - Ae^{ikd} + Be^{-ikd}) &= 2k\gamma\, te^{ikd}.
\end{align}
Using the continuity condition, the derivative condition simplifies to $iBe^{-ikd} = \gamma(Ae^{ikd} + Be^{-ikd})$, giving
\begin{equation}
    B = \frac{\gamma A\, e^{2ikd}}{i - \gamma}. \label{eq:AB2}
\end{equation}

Substituting \eqref{eq:AB2} into \eqref{eq:AB1} and rearranging:
\begin{equation}
    (1+i\gamma)(1 - A) = \gamma A \frac{i - \gamma + \gamma e^{2ikd}}{i - \gamma} \implies A = \frac{1 + i\gamma}{(1+i\gamma)^2 + \gamma^2 e^{2ikd}},
\end{equation}
where we used $(1+i\gamma)^2 = 1 + 2i\gamma - \gamma^2$. The transmission amplitude is $t = A + Be^{-2ikd} = Ai/(i-\gamma) = A/(1+i\gamma)$:
\begin{equation}
    t = \frac{1}{(1 + i\gamma)^2 + \gamma^2 e^{2ikd}}.
\end{equation}

Expanding $(1+i\gamma)^2 = (1-\gamma^2) + 2i\gamma$:
\begin{align}
    \left|(1 + i\gamma)^2 + \gamma^2 e^{2ikd}\right|^2
    &= \big[(1 - \gamma^2) + \gamma^2\cos 2kd\big]^2 + \big[2\gamma + \gamma^2\sin 2kd\big]^2 \nonumber\\
    &= (1 + \gamma^2)^2 + \gamma^4 + 2\gamma^2\big[(1 - \gamma^2)\cos 2kd + 2\gamma\sin 2kd\big].
\end{align}
Define $\theta$ by $\cos 2\theta = -\frac{1-\gamma^2}{1+\gamma^2}$ and $\sin 2\theta = \frac{2\gamma}{1+\gamma^2}$. Then
\begin{equation}
    (1-\gamma^2)\cos 2kd + 2\gamma\sin 2kd = -(1+\gamma^2)\cos 2(kd+\theta),
\end{equation}
and therefore
\begin{equation}
    T_2 = \frac{1}{(1+\gamma^2)^2 + \gamma^4 - 2\gamma^2(1+\gamma^2)\cos 2(kd+\theta)}. \label{eqn:T2-formula}
\end{equation}


\subsection{(1c)}

Define the transfer matrix $M$ by $\binom{A_L}{B_L} = M\binom{A_R}{B_R}$, where the wave on each side is $Ae^{iky} + Be^{-iky}$ with $y$ measured from the scatterer. From (1a), $1 = (1+i\gamma)\,t$ and $r = -i\gamma\, t$, so the transfer matrix for a single delta barrier is
\begin{equation}
    M = \begin{pmatrix} 1+i\gamma & i\gamma \\ -i\gamma & 1-i\gamma \end{pmatrix}.
\end{equation}
The dynamical phase over distance $d$ is encoded in the propagation matrix
\begin{equation}
    U = \begin{pmatrix} e^{-ikd} & 0 \\ 0 & e^{ikd} \end{pmatrix}.
\end{equation}
The total transfer matrix is $M_\mathrm{tot} = M U M$. Computing the $(1,1)$ entry:
\begin{equation}
    [M_\mathrm{tot}]_{11} = (1+i\gamma)^2 e^{-ikd} + \gamma^2 e^{ikd} = e^{-ikd}\big[(1+i\gamma)^2 + \gamma^2 e^{2ikd}\big].
\end{equation}
Since $T_2 = 1/|[M_\mathrm{tot}]_{11}|^2$ and $|e^{-ikd}| = 1$:
\begin{equation}
    T_2 = \frac{1}{\left|(1+i\gamma)^2 + \gamma^2 e^{2ikd}\right|^2},
\end{equation}
which is identical to the result in (1b) before expanding absolute square.


\subsection{(1d)}

Refer to the Equation.~\ref{eqn:T2-formula}, $T_2$ is maximized when $\cos 2(kd + \theta) = 1$, i.e.,
\begin{equation}
    kd + \theta = n\pi, \quad n \in \mathbb{Z}.
\end{equation}
At this maximum:
\begin{equation}
    T_2^{\max} = \frac{1}{(1+\gamma^2)^2 + \gamma^4 - 2\gamma^2(1+\gamma^2)} = \frac{1}{\left[(1+\gamma^2) - \gamma^2\right]^2} = 1.
\end{equation}
Since $T_1 = 1/(1+\gamma^2) < 1$ for any $\gamma \neq 0$, we have $T_2^{\max} > T_1$.

At the resonance condition, the two delta barriers form a cavity, and the particle undergoes multiple reflections between them. When the round-trip phase $2kd + 2\theta$ (dynamical phase $2kd$ and reflection phases $2\theta$) equals a multiple of $2\pi$, the multiply-reflected waves interfere constructively in the forward direction and destructively in the backward direction, yielding perfect transmission---analogous to a Fabry--P\'{e}rot interferometer.


\end{document}
